\documentclass[11pt]{article}
\usepackage[margin=1in]{geometry}
\usepackage{amsmath, amssymb, amsthm}
\usepackage{hyperref}
\usepackage{enumitem}

\newtheorem{theorem}{Theorem}
\newtheorem{proposition}[theorem]{Proposition}
\newtheorem{lemma}[theorem]{Lemma}
\newtheorem{corollary}[theorem]{Corollary}
\newtheorem{conjecture}[theorem]{Conjecture}
\theoremstyle{definition}
\newtheorem{remark}[theorem]{Remark}

\newcommand{\Indep}{\mathrm{Indep}}
\newcommand{\Circ}{\mathrm{Circ}}
\newcommand{\Aut}{\mathrm{Aut}}
\newcommand{\rank}{\mathrm{rank}}
\newcommand{\dualR}{R^{\vee}}
\renewcommand{\H}{\mathcal{H}}
\newcommand{\F}{\mathcal{F}}

\title{Equivariant log-concavity of the f-vector for paving matroids}
\author{}
\date{Draft v3 (2026-05-23)}

\begin{document}
\maketitle

\begin{abstract}
For every paving matroid $M$, we prove that the virtual $\Aut(M)$-representation $[f_{m+1}(M)]^2 - [f_m(M)] \cdot [f_{m+2}(M)]$ is effective for every $m \geq 0$. This establishes \emph{equivariant log-concavity (ELC) of the f-vector for an infinite class of matroids} --- to our knowledge, the first such result beyond direct sums of uniform matroids. The proof reduces ELC to injectivity of a tensor Lefschetz operator $L = \sum x_i \otimes x_i$ acting on a graded vector space $S(M)$, then to a bipartite-incidence rank statement on ordered independent partitions of the ground set. The bipartite-incidence statement is proven via \emph{stressed-hyperplane relaxation} from the uniform case (Karn--Nasr--Proudfoot--Vecchi paradigm) with the boolean Lefschetz theorem as the base case. \S4.5 discusses an attempted descent via the Gui--Xiong equivariant K\"ahler package and explains why it fails: our $L$ is not the descent of the Gui--Xiong operator from the free exterior algebra, so although the symbolic forms match, the equivariant HL on the free exterior algebra does not imply our rank statement on the matroid quotient.
\end{abstract}

A \textbf{paving matroid} is a matroid in which every circuit has size at least equal to the matroid's rank --- equivalently, every subset of size less than the rank is independent. The class includes all \emph{sparse paving} matroids (= paving + co-paving), all Steiner system matroids, and many others; Mayhew--Newman--Welsh--Whittle [MNWW] conjecture that asymptotically almost all matroids are sparse paving (and hence paving).

\section{Introduction}

Mason's f-vector log-concavity conjecture asserts that for every matroid $M$ on $n$ elements, the sequence $f_d(M)$ (= number of independent sets of size $d$) is log-concave:
\[ f_d(M)^2 \geq f_{d-1}(M) \cdot f_{d+1}(M). \]
This was proven by Adiprasito--Huh--Katz (2018) via Hodge theory of the Chow ring; Br\"and\'en--Huh (2020) and Anari--Liu--Oveis Gharan--Vinzant (2018) gave alternative proofs via Lorentzian / completely log-concave polynomials, strengthening to \emph{ultra} log-concavity.

The \textbf{equivariant} lift of Mason's conjecture is the question:

\begin{quote}
\textbf{(ELC for f-vector).} For every matroid $M$ and every $m \geq 0$, is the virtual $\Aut(M)$-representation $[f_{m+1}(M)]^2 - [f_m(M)] \cdot [f_{m+2}(M)]$ effective (i.e., a genuine representation)?
\end{quote}

Here $[f_d(M)] := k[\Indep_d(M)]$ denotes the permutation representation of $\Aut(M)$ on size-$d$ independent sets.

Equivariant log-concavity (ELC) has been established for various \emph{adjacent} invariants --- Orlik--Solomon, Orlik--Terao, and Cordovil algebras of the braid arrangement [MMPR]; equivariant Kazhdan--Lusztig polynomials of paving matroids [KNPV]; equivariant Kazhdan--Lusztig polynomials of q-niform and uniform matroids [GLXYZ] --- but the f-vector statement specifically has remained open except in three trivial cases:

\begin{enumerate}[leftmargin=*]
\item Boolean matroids $U_{n,n}$ (= power set of $[n]$).
\item Uniform matroids $U_{r,n}$, by inheritance from the boolean case.
\item Direct sums of uniforms, by K\"unneth.
\end{enumerate}

The main result of this note establishes ELC of the f-vector for a much larger class:

\begin{theorem}
Let $M$ be a \textbf{paving matroid}: a matroid in which every circuit has size at least $\rank(M)$. Equivalently, every subset of size less than $\rank(M)$ is independent. Then for every $m \geq 0$, the virtual $\Aut(M)$-representation $[f_{m+1}(M)]^2 - [f_m(M)] \cdot [f_{m+2}(M)]$ is effective.
\end{theorem}

Paving matroids are a substantial class. They include all sparse paving matroids (= paving with paving dual), which by [MNWW] are conjecturally asymptotically almost all matroids. They also include Steiner system matroids, and many specific designs.

The proof proceeds in three steps:
\begin{enumerate}[leftmargin=*]
\item \S2: Identify ELC as the injectivity of a single explicit $\Aut(M)$-equivariant linear operator $L$ on a graded vector space $S(M)$.
\item \S3: Decompose $L$ into orbits indexed by pairs $(C, U)$ and reduce each orbit's injectivity to a bipartite-incidence rank statement on ``balanced indep bipartitions'' of the contracted-restricted minor.
\item \S4: Prove the bipartite-incidence statement for paving matroids via a single inductive step from the uniform matroid, in which the entire collection of circuit-hyperplanes is added at once. The boolean Lefschetz theorem supplies the base case; a submatrix-rank argument supplies the inductive step.
\end{enumerate}

\section{The Lefschetz operator and the ELC criterion}

Let $M$ be a loopless matroid on $E = \{1, \dots, n\}$ with rank $r$. Denote by $\Indep_d(M)$ the set of size-$d$ independent sets and $f_d(M) := |\Indep_d(M)|$. The $\Aut(M)$-permutation representation on $\Indep_d(M)$ is $[f_d(M)]$.

\subsection{The ring $R(M)$ and its dual}

Set
\[ R(M) := k[x_1, \dots, x_n] \big/ \big(x_C : C \in \Circ(M);\ x_i^2 : i \in E\big), \]
where $\Circ(M)$ is the set of circuits and $x_S := \prod_{i \in S} x_i$. The squarefree quotient ensures $R(M)_d$ has basis $\{x_S : S \in \Indep_d(M)\}$, so $\dim R(M)_d = f_d(M)$. As $\Aut(M)$-representations, $R(M)_d \cong [f_d(M)]$.

Let $\dualR(M)$ be the graded linear dual, with basis $\{y_S : S \in \Indep_d(M)\}$ dual to $\{x_S\}$. The matroid algebra $R(M)$ acts on $\dualR(M)$ by contraction: $x_i \cdot y_S = y_{S \setminus \{i\}}$ if $i \in S$, else $0$.

\subsection{$S(M)$ and the Lefschetz operator}

Define $S(M) := R(M) \otimes_k \dualR(M)$, bigraded by \textbf{external degree} $e := |S| - |T|$ and \textbf{internal degree} $m := |S|$ on $x_S \otimes y_T$. The \textbf{Lefschetz operator} is
\[ L := \sum_{i=1}^n x_i \otimes x_i \in \mathrm{End}_k(S(M)), \]
with explicit action
\[ L(x_S \otimes y_T) = \sum_{i \in T \setminus S,\ S \cup \{i\} \in \Indep(M)} x_{S \cup \{i\}} \otimes y_{T \setminus \{i\}}. \]
$L$ raises external degree by 2, raises internal degree by 1, and is $\Aut(M)$-equivariant.

\subsection{ELC $\iff$ Injectivity of $L$}

For each $m, d \geq 0$ with $d \geq 2$, consider the block
\[ L_m : R(M)_m \otimes \dualR(M)_{-(m+d)} \to R(M)_{m+1} \otimes \dualR(M)_{-(m+d-1)}. \]
The source is $[f_m(M)] \cdot [f_{m+d}(M)]^* \cong [f_m(M)] \cdot [f_{m+d}(M)]$ (the dual of a permutation representation equals itself); the target is $[f_{m+1}(M)] \cdot [f_{m+d-1}(M)]$.

\begin{proposition}
ELC for the f-vector of $M$ holds if and only if, for every $m \geq 0$ and every $d \geq 2$, $L_m$ is injective. Taking $d = 2$ gives the standard log-concavity form: cokernel of $L_m$ equals $[f_{m+1}(M)]^2 - [f_m(M)] \cdot [f_{m+2}(M)]$, effective iff $L_m$ injective.
\end{proposition}

\begin{proof}
If $L_m$ is $\Aut(M)$-equivariant and injective, its cokernel is an honest $\Aut(M)$-representation (= effective virtual). The converse is standard.
\end{proof}

This reduces ELC to the injectivity of $L$ on each bigrade $(m, -d)$.

\section{Orbit decomposition and the X-bipartite incidence}

\subsection{The orbit structure}

For $(S, T)$ a basis element of $R(M)_m \otimes \dualR(M)_{-(m+d)}$ (so $|S| = m$, $|T| = m+d$, both independent), set $C := S \cap T$ and $U := S \cup T$.

\begin{lemma}
$L$ preserves the pair $(C, U)$: if $(S, T) \mapsto (S \cup \{i\}, T \setminus \{i\})$ is a summand of $L(x_S \otimes y_T)$, then $(S \cup \{i\}) \cap (T \setminus \{i\}) = C$ and $(S \cup \{i\}) \cup (T \setminus \{i\}) = U$.
\end{lemma}

\begin{proof}
The summand requires $i \in T \setminus S$. Then $S' \cup T' = U$ since $i \in T$, and $S' \cap T' = (C \cup \{i\}) \setminus \{i\} = C$ since $i \notin C$.
\end{proof}

Define $N := N(C, U) := (M / C) \big|_{U \setminus C}$, the contracted-restricted matroid on ground set $U \setminus C$ with $\Indep(N) = \{A \subseteq U \setminus C : A \cup C \in \Indep(M)\}$.

Set $k := m - |C|$ and $n' := |U \setminus C| = 2k + d$. The orbit $O_{C, U}$ is in bijection (via $(S, T) \mapsto A := S \setminus C \subseteq U \setminus C$) with the set of \textbf{balanced indep bipartitions}
\[ X_k(N) := \big\{A \subseteq U \setminus C : |A| = k,\ A \in \Indep(N),\ (U \setminus C) \setminus A \in \Indep(N)\big\}. \]

\subsection{Reduction to the X-bipartite incidence}

\begin{proposition}
Under the bijection above, $L|_{O_{C, U}}$ corresponds to the operator
\[ \partial^* : \mathbb{R}^{X_k(N)} \to \mathbb{R}^{X_{k+1}(N)}, \quad \partial^*(x_A) = \sum_{i \in (U \setminus C) \setminus A,\ A \cup \{i\} \in \Indep(N)} x_{A \cup \{i\}}. \]
Moreover, $\partial^*(\mathbb{R}^{X_k(N)}) \subseteq \mathbb{R}^{X_{k+1}(N)}$ (the image automatically lands in the X-restricted subspace at the next level).
\end{proposition}

\begin{proof}
The summands of $L|_{O_{C, U}}$ correspond to moving an element $i \in T \setminus S = (U \setminus C) \setminus A$ from $T$ to $S$, with the matroid constraint $S \cup \{i\} \in \Indep(M) \Leftrightarrow A \cup \{i\} \in \Indep(N)$. The complement condition for the result to remain in $X_{k+1}(N)$ is automatic: $(U \setminus C) \setminus (A \cup \{i\}) \subseteq (U \setminus C) \setminus A \in \Indep(N)$, so it is independent.
\end{proof}

In view of the above:

\begin{theorem}[Reduction]
ELC for the f-vector of $M$ holds if and only if, for every minor $N = N(C, U)$ of $M$ arising from an orbit pair $(C, U)$ with $2k + d = |E(N)|$, $k + d \leq \rank(N)$, $d \geq 2$, the operator $\partial^* : \mathbb{R}^{X_k(N)} \to \mathbb{R}^{X_{k+1}(N)}$ is injective.
\end{theorem}

The remaining work is to prove the injectivity statement. Minors of paving matroids are themselves paving, so it suffices to prove:

\begin{theorem}
For every paving matroid $N$ on $n' = 2k + d$ elements of rank $r' \geq k + d$, with $d \geq 2$, the bipartite-incidence operator $\partial^* : \mathbb{R}^{X_k(N)} \to \mathbb{R}^{X_{k+1}(N)}$ is injective.
\end{theorem}

\section{Stressed-hyperplane relaxation for paving matroids}

A paving matroid $N$ with rank $r'$ on $n'$ elements can be described by its set $\H$ of \textbf{circuit-hyperplanes} (CHs): the $r'$-subsets of $E(N)$ that are \emph{not} bases (= dependent of full rank). Since $N$ is paving, every subset of size $< r'$ is independent, and the dependent $r'$-subsets are exactly the CHs in $\H$. We write $N = U_{r', n'}(\H)$.

\begin{remark}[On validity]
Not every collection $\H$ of $r'$-subsets gives a valid matroid. For two CHs $F_1, F_2 \in \H$ with $|F_1 \cap F_2| = r' - 1$, the matroid axioms force additional CHs. \textbf{Sparse paving} is the special case where all pairwise CH intersections satisfy $|F_1 \cap F_2| \leq r' - 2$, so CHs can be added one at a time. Our proof handles both cases uniformly.
\end{remark}

\subsection{The base case: uniform matroids}

\begin{lemma}
If $N = U_{r', n'}$ (so $\H = \emptyset$), then Theorem 6 holds: $\partial^* : \mathbb{R}^{X_k(N)} \to \mathbb{R}^{X_{k+1}(N)}$ is injective.
\end{lemma}

\begin{proof}
For the uniform matroid $U_{r', n'}$, every subset of size $\leq r'$ is independent. Since $n' = 2k + d \leq 2r'$, both $|A| = k \leq r'$ and $|E \setminus A| = k + d \leq r'$ are satisfied. Hence $X_k(U_{r', n'}) = \binom{E(N)}{k}$ consists of all $k$-subsets, and similarly $X_{k+1}$.

The bipartite-incidence operator on these spaces is exactly the simplicial coboundary on the boolean lattice at level $k \to k+1$, which is injective for $k < n'/2$. Since $n' = 2k + d$ with $d \geq 2$, indeed $k < n'/2$. So $\partial^*$ is injective by the standard boolean $\mathfrak{sl}_2$-Lefschetz theorem.
\end{proof}

\subsection{Inductive step: adding stressed hyperplanes}

\begin{lemma}
Let $N = U_{r', n'}(\H)$ be a paving matroid, and let $N' = U_{r', n'}(\H \cup \F)$ be another paving matroid obtained by adding a nonempty collection $\F$ of new CHs (each $F \in \F$ is currently a basis of $N$). Suppose Theorem 6 holds for $N$. Then Theorem 6 holds for $N'$.
\end{lemma}

\begin{proof}
Fix $(k, d)$ with $2k + d = n'$, $r' \geq k + d$, $d \geq 2$.

\textbf{Claim (a):} $X_{k+1}(N') = X_{k+1}(N)$.

Going from $N$ to $N'$, the only sets that lose independent status are those in $\F$. A set $A' \in X_{k+1}(N)$ is removed iff $A' \in \F$ or $E(N) \setminus A' \in \F$. Now $|A'| = k + 1$ and each $F \in \F$ has size $r' \geq k + 2 > k + 1$, so $A' \notin \F$. Also $|E(N) \setminus A'| = k + d - 1 < r'$, so $E(N) \setminus A' \notin \F$.

\textbf{Claim (b):} $|X_k(N')| = |X_k(N)| - \delta$ where $0 \leq \delta \leq |\F|$. The removed elements are exactly $\{E(N) \setminus F : F \in \F, E(N) \setminus F \in X_k(N)\}$, non-empty only in the boundary case $r' = k + d$.

A set $A \in X_k(N)$ is removed iff $A \in \F$ or $E(N) \setminus A \in \F$. Since $|A| = k < r'$, we have $A \notin \F$. The condition $E(N) \setminus A \in \F$ requires $|E(N) \setminus A| = r'$, i.e., $k = n' - r'$. By hypothesis $r' \geq k + d$, so this has equality iff $r' = k + d$.

\textbf{Inductive step proper.} By Claims (a)--(b), the bipartite-incidence matrix $\partial^* : \mathbb{R}^{X_k(N')} \to \mathbb{R}^{X_{k+1}(N')}$ is a submatrix of the bipartite-incidence matrix for $N$, obtained by removing $\delta$ columns and no rows. The original matrix has full column rank by hypothesis, so its columns are linearly independent. Any subset of these columns is also linearly independent, so the new matrix has full column rank, i.e., $\partial^*$ injective.
\end{proof}

\subsection{Proof of Theorem 1}

Let $N = U_{r', n'}(\H)$ be any paving matroid with the bigrade hypothesis. Take $\F := \H$ and apply Lemma 8 in a single step from $U_{r', n'}$: the base case (Lemma 7) and one inductive step (Lemma 8) give Theorem 6 for $N$. Combined with Theorem 5 (the reduction) and the fact that minors of paving matroids are paving, this gives ELC of the f-vector for every paving matroid $M$.

\subsection{Non-sparse-paving example (verification)}

To illustrate the proof beyond sparse paving, consider $M = U_{5, 8}(\H)$ where $\H$ consists of all six 5-subsets of $\{0, 1, 2, 3, 4, 5\}$. Pairwise these CHs intersect in 4 elements, violating the sparse paving condition. $M$ cannot be built by adding the CHs one at a time and getting valid intermediate matroids; but adding all six together gives a valid paving matroid, and Lemma 8's submatrix-rank argument applies in one inductive step.

Computational verification (in \texttt{computations/test\_paving\_extension.py}): at bigrade $(k, d) = (3, 2)$, $|X_3(M)| = 50$, $|X_4(M)| = 70$, and the bipartite-incidence rank is exactly 50 = $|X_3(M)|$.

\subsection{An attempted descent via Gui--Xiong's framework, and why it fails}

It is natural to ask whether Theorem 6 follows from the Gui--Xiong equivariant K\"ahler package [GX, arXiv:2205.05420], which gives Hard Lefschetz on the free exterior algebra $\Lambda(V) \otimes \Lambda(V^*)$ with the operator $L_{\mathrm{amb}} := \sum e_{\theta_k} \otimes i_{\theta_k}$, equivariant under any group acting on $V$. Symbolically, this $L_{\mathrm{amb}}$ is the same expression as our $L$, so one might hope that our $L$ on $R(M) \otimes \dualR(M)$ is the descent of $L_{\mathrm{amb}}$ under the projection $\pi: \Lambda(V) \otimes \Lambda(V^*) \to R(M) \otimes \dualR(M)$.

\textbf{This fails.} Our $L$ on $R(M) \otimes \dualR(M)$ (defined intrinsically as the unsigned bipartite incidence on indep-set pairs) is \emph{not} the descent of $L_{\mathrm{amb}}$. Concretely: take $M = U_{5,8}$ with circuit-hyperplane $F = \{0,1,2,3,4\}$ and $v = x_{\{5,6,7\}} \otimes y_F$. Then $v \in \ker(\pi)$ (since $y_F = 0$ in $\dualR(M)$), but
\[ L_{\mathrm{amb}}(v) = \sum_{i \in F} x_{\{5,6,7\} \cup \{i\}} \otimes y_{F \setminus \{i\}}, \]
and every summand has both factors \emph{independent} in $M$. Hence $L_{\mathrm{amb}}(v) \notin \ker(\pi)$, so the diagram doesn't commute. Verification in \texttt{computations/verify\_descent.py}.

\section{Discussion}

\subsection{Relationship with prior work}

The technique of stressed-hyperplane relaxation was developed by Karn--Nasr--Proudfoot--Vecchi [KNPV] for proving equivariant log-concavity of equivariant Kazhdan--Lusztig polynomials of paving matroids. Our application to the f-vector is, to our knowledge, the first time this technique has been applied to the permutation representation $[f_d(M)]$.

The dimensional (= non-equivariant) part of our main theorem was recently established by Ardila-Mantilla--Cristancho--Denham--Eur--Huh--Wang [ACDEHW] via Lorentzian polynomial methods. Their result handles dimensional inequality for \emph{all} matroids; the additional rank statement (= injectivity of $\partial^*$) is what we contribute for the paving class.

\subsection{The X-restricted rank problem (open)}

The cleanest open problem to attack is the rank statement directly, for non-paving matroids.

\begin{conjecture}
For every loopless matroid $M$ and every bigrade $(m, -d)$ with $2m + d \leq n$, $m + d \leq r$, $d \geq 2$, the bipartite-incidence operator $\partial^*: \mathbb{R}^{X_m(M)} \to \mathbb{R}^{X_{m+1}(M)}$ is injective.
\end{conjecture}

Computational evidence: across more than $2 \times 10^5$ matroid orbits checked --- including graphic, projective, V\'amos, Fano, Pappus, NonPappus, and various restrictions of $M(K_6)$ --- no counterexample has emerged.

\subsection{Open problems}

\begin{enumerate}[leftmargin=*]
\item Generalize to arbitrary matroids. A Schur-complement formulation may be needed.
\item Find a direct geometric realization of the X-bipartite incidence operator as a Lefschetz-type morphism, \`a la AHK or Eur--Huh--Larson.
\item Quantify the gap: for non-paving matroids, even though the relaxation argument fails, the rank deficit of the bipartite incidence on $\Indep_k(M) \to \Indep_{k+1}(M)$ can be substantial, while the X-restricted operator empirically still has full rank. Understand the structural reason for the ``X-restoration''.
\end{enumerate}

\begin{thebibliography}{99}

\bibitem[AHK]{AHK} K. Adiprasito, J. Huh, E. Katz. \emph{Hodge theory for combinatorial geometries.} Annals of Mathematics 188 (2018), 381--452.

\bibitem[ALOGV]{ALOGV} N. Anari, K. Liu, S. Oveis Gharan, C. Vinzant. \emph{Log-Concave Polynomials, III: Mason's Ultra-Log-Concavity Conjecture for Independent Sets of Matroids.} arXiv:1811.01600.

\bibitem[ACDEHW]{ACDEHW} F. Ardila-Mantilla, S. Cristancho, G. Denham, C. Eur, J. Huh, B. Wang. \emph{Tree metrics and log-concavity for matroids.} arXiv:2601.02547 (January 2026).

\bibitem[BH]{BH} P. Br\"and\'en, J. Huh. \emph{Lorentzian polynomials.} Annals of Mathematics 192 (2020), 821--891.

\bibitem[GLXYZ]{GLXYZ} A. L. L. Gao, E. Y. H. Li, M. H. Y. Xie, A. L. B. Yang, Z.-X. Zhang. \emph{Induced log-concavity of equivariant matroid invariants.} arXiv:2307.10539.

\bibitem[GX]{GX} L. Gui, R. Xiong. \emph{Equivariant log-concavity and equivariant K\"ahler packages.} arXiv:2205.05420.

\bibitem[Sw]{Sw} E. Swartz. \emph{g-elements of matroid complexes.} arXiv:math/0210376.

\bibitem[ANR]{ANR} N. Angarone, A. Nathanson, V. Reiner. \emph{Chow rings of matroids as permutation representations.} arXiv:2309.14312.

\bibitem[KNPV]{KNPV} T. Karn, G. Nasr, N. Proudfoot, L. Vecchi. \emph{Equivariant Kazhdan--Lusztig theory of paving matroids.} arXiv:2202.06938.

\bibitem[MMPR]{MMPR} J. Matherne, D. Miyata, N. Proudfoot, E. Ramos. \emph{Equivariant log concavity and representation stability.} IMRN 2023, 3885--3906.

\bibitem[MNWW]{MNWW} D. Mayhew, M. Newman, D. Welsh, G. Whittle. \emph{On the asymptotic proportion of connected matroids.} European Journal of Combinatorics 32 (2011), 882--890.

\end{thebibliography}

\end{document}
